\section{Contact lines and a suitable pair}
\label{lines:section}

Throughout this section assume~\eqref{red:primitive}. A
\emph{contact line} is a map $f:\PP^1\to X$ with $f^*L=\OO(1)$,
as in Section~\ref{sec:introduction}.
Every such map is birational onto its image, since a nontrivial
cover would multiply the positive integral $L$-degree. It is
horizontal:
\begin{equation}\label{lines:horizontal}
\theta\circ df=0,
\end{equation}
because $\operatorname{Hom}(\OO(2),\OO(1))=0$.

Our goal is the pair in Proposition~\ref{lines:suitable-pair} at a general
point $x$: two free contact lines whose positive tangent-bundle subspaces
span $F_x$ and whose tangent vectors have nonzero Levi pairing.

We first specify the published contact-line results used below.
Kebekus~\cite[Section~2.3]{Kebekus2001} constructs an irreducible covering
component $\mathcal H$ of rational curves on a projective contact
Fano manifold with $b_2=1$. Its minimal contact degree is one or two;
degree two forces projective space. The latter has contact line bundle
$\OO(2)$ and is excluded by~\eqref{red:primitive}. Thus we may fix
a covering component of contact lines. The parameter space of
these curves through a general point is proper. For a general
$x\in X$, Kebekus's Proposition~3.3 gives smoothness of the lines
through $x$, Corollary~3.4 supplies their tangent-direction morphism,
and Lemma~3.5 gives the splitting
\begin{equation}\label{lines:standard-splitting}
f^*T_X\simeq
\OO(2)\oplus\OO(1)^{\oplus(n-1)}\oplus\OO^{\oplus(n+1)}.
\end{equation}
Theorem~4.4 of \cite{Kebekus2001} says that the tangent directions
of the curves from $\mathcal H$ through $x$ span $\PP(F_x)$.
This is a statement about their union; no individual irreducible
component of the set of directions is asserted to span.

For $n\geq2$ we also use the corrected Theorem~3.1 of
Buczy\'nski~\cite{Buczynski2010}, in the consolidated version
arXiv:1002.0698v2. It says that the reduced variety $Z_x$ of these
tangent directions at a general point is smooth and Legendrian
in the contact projective space $\PP(F_x)$, and that its points
are actual tangents to contact lines through $x$. The result
allows a union of components of the family of contact lines.
We use neither irreducibility of $Z_x$ nor an isomorphism between
the line locus and its tangent cone; see the correction discussed
in Remark~3.2 of that paper. Smoothness means that the irreducible
components of $Z_x$ are disjoint. Each has dimension $n-1$, and
its affine cone is smooth away from the vertex, with Lagrangian
tangent spaces for the symplectic form on $F_x$.

We record the deformation details needed to use these directions.
For a morphism from a smooth projective curve to a smooth variety,
the tangent space of the map scheme is $H^0(f^*T_X)$; with specified
point values it is the subspace of sections vanishing there.
The corresponding $H^1$ is an obstruction space. In particular
$H^1(f^*T_X)=0$ implies smoothness of the map scheme at $f$;
see~\cite[Theorem~5]{AraujoKollar2003}. These statements follow
from local liftings across a square-zero extension: differences
of liftings form a \v{C}ech cocycle in the pulled-back tangent
bundle, and its $H^1$ class is the obstruction to gluing.

\begin{lemma}[Freeness in fixed degree]\label{lines:fixed-degree}
Let $A$ be an ample line bundle on a smooth projective variety $M$,
and fix an integer $e>0$. There is a dense open $M_e\subset M$
such that every map $f:\PP^1\to M$ of $A$-degree $e$ whose image
meets $M_e$ is free, that is, $f^*T_M$ is globally generated.
\end{lemma}

\begin{proof}
The degree-$e$ map scheme is of finite type: an embedding by a
very ample power $A^a$ realizes it as the locus of maps into $M$
among degree-$ae$ maps to a fixed projective space. Stratify its
reduction into finitely many smooth irreducible locally closed
sets $S_j$. Consider evaluation at $0$ on each stratum.
Discard the closures of the images of nondominant strata.
For a dominant stratum discard also the closure of the image of
its critical locus. The latter closure is proper: otherwise
some reduced irreducible component $C$ of that critical locus
would dominate $M$. On a smooth dense open of $C$, generic
smoothness in characteristic zero would give a surjective
differential to $M$. It factors through the differential on
$S_j$, contradicting criticality.

The complement $M_e$ of these finitely many proper closed sets
is dense and open. At any map with $f(0)\in M_e$, stratum
tangent vectors give global sections of $f^*T_M$ whose values
span $T_{f(0)}M$. In the splitting
$f^*T_M=\bigoplus_i\OO(a_i)$, a negative summand would have no
global sections and could not be generated at $0$. Hence every
$a_i\geq0$, giving global generation. Reparametrizing $\PP^1$
moves any specified point to $0$, proving the assertion for
every map whose image meets $M_e$.
\end{proof}

\begin{lemma}\label{lines:positive}
For a free contact line $f$, let $P_f\subset f^*T_X$ be the
positive part of its splitting. This subbundle is intrinsic,
the quotient $f^*T_X/P_f$ is trivial, and
\begin{equation}\label{lines:positive-horizontal}
\theta(P_f)=0.
\end{equation}
Moreover $P_f$ is isotropic for the pulled-back Levi form, so
$\operatorname{rank}P_f\leq n$.
\end{lemma}

\begin{proof}
Freeness says that all splitting degrees are nonnegative.
The sum of the positive summands is independent of the splitting:
any homomorphism from a positive line bundle to a trivial bundle
is zero, so the positive part of one splitting is contained in
that of any other. The quotient is the sum of the degree-zero
summands and is trivial.

The nonzero differential $df:\OO(2)\to f^*F$, paired with the
nondegenerate Levi form, gives a nonzero homomorphism
\[
f^*F\longrightarrow
T^*_{\PP^1}\otimes f^*L=\OO(-1).
\]
Consequently $f^*F$ is not globally generated: all global sections
would map to zero in $\OO(-1)$. If $\theta$ were nonzero on a
positive summand $\OO(a)$, then $a=1$ and that map to $f^*L=\OO(1)$
would be an isomorphism. It would split the pulled-back contact
sequence. The kernel $f^*F$ would then be a direct summand of the
globally generated bundle $f^*T_X$, a contradiction. This proves
\eqref{lines:positive-horizontal}. Finally each component of the
Levi pairing on positive summands is a homomorphism
$\OO(a+b)\to\OO(1)$ with $a,b\geq1$, and is zero. Isotropy and
nondegeneracy of the rank-$2n$ Levi form give the rank bound.
\end{proof}

\begin{corollary}\label{lines:free-cover}
There is a smooth irreducible family $H$ of parametrized free
contact lines with smooth dominant evaluation
$H\times\PP^1\to X$.
\end{corollary}

\begin{proof}
Take a component of the map scheme corresponding to the covering
family $\mathcal H$. Lemma~\ref{lines:fixed-degree} makes its free
locus nonempty. Freeness is open, and at a free map both
$H^1(f^*T_X)$ and $H^1(f^*T_X(-t))$ vanish for every $t\in\PP^1$.
The first vanishing makes this locus smooth; the second makes
evaluation smooth at every pair $(f,t)$, by the deformation
criterion above. A nonempty irreducible open in this locus gives
$H$, and its evaluation remains dominant.
\end{proof}

\subsection{The dimension-three endpoint}

When $n=1$, the needed description of directions follows directly
from Kebekus's results, which have no $n\geq2$ hypothesis.
Equation~\eqref{lines:standard-splitting} becomes
$f^*T_X=\OO(2)\oplus\OO^{\oplus2}$. For a general $x$ and a line
parametrization with $f(0)=x$, the space of maps with this value
fixed is smooth of dimension
\[
h^0(\PP^1,f^*T_X(-0))=h^0(\PP^1,\OO(1))=2.
\]
The reparametrizations fixing $0$ form a two-dimensional group.
They act freely on each parametrization, since it is birational
onto its image. Their orbits are consequently open in the smooth
fixed-value map space near these maps. The projective tangent
direction is constant on each orbit. Equivalently the local
space of line images through $x$ has dimension zero; the proper
parameter space of such images is therefore finite. Its reduced
tangent-direction image $Z_x\subset\PP(F_x)=\PP^1$ is a finite
smooth set.
Kebekus's spanning Theorem supplies two distinct directions.
Their affine cones are two distinct lines in the symplectic
plane $F_x$, hence have nonzero mutual symplectic pairing and
sum equal to $F_x$. Each is Lagrangian.
This gives the required cone geometry in dimension three without
using the higher-dimensional smoothness assertion.

\subsection{Addition of two Legendrian cones}

The secant geometry behind this argument appears in
Landsberg--Manivel~\cite[Proposition~17(2) and its proof]{LandsbergManivel2007}
and Buczy\'nski~\cite[Proposition~3.5]{Buczynski2010}.
We prove the pairwise assertion below for a possibly disconnected
Legendrian variety; it will apply to the entire direction variety
$Z_x$ without requiring any one of its components to span.

\begin{lemma}\label{lines:cone-addition}
Let $(V,\omega)$ be a symplectic vector space of dimension $2n$.
Let $Z\subset\PP(V)$ be a smooth projective Legendrian variety,
possibly disconnected. Write $U_i$ for the affine cones over
its irreducible components. If $\omega(a,b)\neq0$ for some
$a\in U_i$, $b\in U_j$, then the addition map
\[
U_i\times U_j\longrightarrow V,\qquad (a,b)\longmapsto a+b
\]
is dominant.
\end{lemma}

\begin{proof}
Each cone has dimension $n$, is smooth away from its vertex,
and has Lagrangian tangent spaces there. Suppose addition is
not dominant. Its source has dimension $2n$, so its general
fibers have positive dimension. The nonorthogonal pairs form a
nonempty open subset of the irreducible source. Choose a general such
pair in a fiber of positive dimension, and an irreducible
algebraic curve $C$ in that fiber through the pair. Such a curve
is obtained by intersecting an affine neighborhood in a
positive-dimensional fiber component with sufficiently many
general hyperplanes through the point. Write its two coordinate
maps as $a(t),b(t)$, with fixed sum $s\in V$.

On the smooth dense open of $C$ where both coordinates are
nonzero, their derivatives satisfy $\dot a=-\dot b$.
The radial vector of a cone belongs to its tangent space.
Lagrangianity therefore implies
\[
\frac{d}{dt}\omega(a,b)
=\omega(\dot a,b)+\omega(a,\dot b)=0:
\quad \omega(\dot a,b)=-\omega(\dot b,b)=0,
\quad \omega(a,\dot b)=-\omega(a,\dot a)=0.
\]
In characteristic zero a rational function on an integral curve
with zero differential is constant. Hence $\omega(a,b)=c\neq0$
on $C$. In particular neither coordinate is ever zero there.

Normalize $C$ and take its smooth projective completion
$\overline C$. The vector-valued rational function $a$ is
nonconstant, since $a+b=s$ and $C$ is a curve of pairs. At least
one coordinate of $a$ has a pole on $\overline C$, as regular
functions on a connected projective curve are constant. Choose
such a place, a uniformizer $t$, and let $k>0$ be the maximal
pole order of the coordinates of $a$. Then
\[
p(t)=t^ka(t),\qquad q(t)=-t^kb(t)
\]
are regular near $0$, with
\begin{equation}\label{lines:rescaling}
p(0)=q(0)\neq0,\qquad p(t)-q(t)=t^ks.
\end{equation}
The closed conical property puts $p(0)$ in $U_i$ and $q(0)$ in
$U_j$. Distinct components of the smooth variety $Z$ are
disjoint. Their common nonzero point therefore forces $i=j$.
Denote the resulting smooth cone germ by $U$.

We claim that $\omega(p,q)$ on $U\times U$ vanishes to order at
least three along the diagonal. In local coordinates write
$r(z)$ for the inclusion of $U$ into $V$. Conical
Lagrangianity gives
\[
\omega(r,r_i)=0,\qquad \omega(r_j,r_i)=0,
\]
where subscripts denote coordinate derivatives. Differentiating
the first identity gives $\omega(r,r_{ij})=0$. Taylor expansion
of $\omega(r(z),r(z+h))$ has zero terms of degrees zero, one,
and two in $h$. These identities hold for every $z$, so the
function belongs to the cube of the diagonal ideal on $U\times U$;
this proves the claim even when both arguments vary.
Choose the local coordinates $z$ to be restrictions of linear
coordinates on $V$ with
independent differentials on $T_{p(0)}U$. Equation~\eqref{lines:rescaling}
then says that the local coordinates of $p(t)$ and $q(t)$ differ
by $O(t^k)$. The claim yields
\[
\omega(p(t),q(t))=O(t^{3k}).
\]
But its exact value is $-t^{2k}\omega(a(t),b(t))=-ct^{2k}$,
of order $2k$. This contradiction proves dominance.
\end{proof}

\begin{proposition}[Suitable contact-line pair]\label{lines:suitable-pair}
There exist a point $x\in X$ and free parametrized contact lines
$f_1,f_2:\PP^1\to X$, with $f_1(0)=f_2(0)=x$, such that, for
affine parameters at $0$ and $a=f'_1(0)$, $b=f'_2(0)$,
\begin{equation}\label{lines:pair}
d\alpha_x(a,b)\neq0,\qquad
P_{f_1}|_0+P_{f_2}|_0=F_x.
\end{equation}
In particular the two positive subspaces are complementary
Lagrangian subspaces of $F_x$.
\end{proposition}

\begin{proof}
Choose $x$ in the common general locus of the contact-line results
and Lemma~\ref{lines:fixed-degree} for degree one, so that every
line map through $x$ is free. Fix a local frame of $L$ and put
$\omega=d\alpha_x|_{F_x}$. The description of $Z_x$ above applies
for $n\geq2$, and its direct dimension-three replacement applies
for $n=1$. Since $Z_x$ spans $\PP(F_x)$, some pair of its cone
components $U_i,U_j$ has $\omega(a,b)\neq0$ for a pair of points:
otherwise bilinearity would make $\omega$ zero on all of $F_x$.
Lemma~\ref{lines:cone-addition} makes their addition dominant
(for $n=1$ it is already a linear isomorphism).

Consider the finite-type reduced space of maps with $f(0)=x$.
Taking the derivative at $0$ defines an algebraic map to $F_x$.
Every nonzero vector in $U_i$ or $U_j$ is realized: each projective
direction comes from an actual smooth contact line, and
reparametrizations fixing $0$ multiply its derivative by any
nonzero scalar. Thus we can choose irreducible locally closed
subspaces of the map space whose derivative images are dense
in $U_i$ and $U_j$. On their product the derivative-sum map
is dominant. Restricting to the nonempty open where the two
derivatives have nonzero pairing preserves dominance.
Generic smoothness on a smooth dense open of this reduced
product supplies a pair where its differential surjects onto
$F_x$.

It remains to identify where that differential takes values.
A tangent vector to maps with $f(0)=x$ is a section
$\psi\in H^0(\PP^1,f^*T_X(-0))$. Its projection to the trivial
quotient $f^*T_X/P_f$ is a constant section vanishing at $0$,
hence zero everywhere. Thus $\psi$ takes values in $P_f$.
Because $\psi(0)=0$, its derivative at $0$ is intrinsically in
$P_f|_0$: differentiating a local frame of $P_f$ contributes
only terms multiplied by the zero value of $\psi$. In target
coordinates this derivative is exactly the variation of
$f'(0)$. Consequently the surjective differential of the
sum map has image contained in $P_{f_1}|_0+P_{f_2}|_0$.
Lemma~\ref{lines:positive} puts this sum inside $F_x$, proving
\eqref{lines:pair}. Each summand is isotropic of dimension at
most $n$; since their sum has dimension $2n$, both have
dimension $n$ and their intersection is zero.
\end{proof}
